\documentclass[aps,pre,reprint,nofootinbib,floatfix]{revtex4-2}

\usepackage{amsmath,amssymb,bm}
\usepackage{array}
\usepackage{booktabs}
\newcolumntype{L}[1]{>{\raggedright\arraybackslash}p{#1}}
\usepackage{graphicx}
\usepackage{hyperref}
\usepackage{microtype}
\usepackage{pgfplots}
\usepgfplotslibrary{groupplots}
\pgfplotsset{compat=1.15,
  paperaxis/.style={
    width=\columnwidth, height=0.58\columnwidth,
    grid=major, grid style={gray!20},
    tick label style={font=\footnotesize},
    label style={font=\footnotesize},
    legend style={font=\footnotesize, draw=gray!45, fill=white},
    line width=0.9pt}}

\raggedbottom

\begin{document}

\title{What a Projection Fold Can and Cannot Imitate\\
The Complete Arc of a Campaign's Founding Question, from a Kinematic
Theorem to a Sealed Analyticity Exponent}

\author{Andrew H. Bond}
\email{andrew.bond@sjsu.edu}
\thanks{ORCID: 0009-0003-2599-6158}
\affiliation{San Jos\'e State University, San Jos\'e, California, USA}
\date{August 7, 2026}

\begin{abstract}
A smooth worldline in a hidden manifold, observed through a projection
that folds with respect to a detector's time, produces two apparent
branches of opposite orientation where there was one.  That is a
theorem about maps.  Whether any local hidden dynamics makes such folds
at a rate with the structure physical pair production has is not a
theorem, and this paper reports the whole attempt to find out, across
twenty-nine runs on declared finite models, with every failure kept.  The
instrument net reproduces every exact control, a
degenerate critical point is never classified as a fold, and a thousand
perturbed members of the negative control each produce exactly twelve
folds, which is what a rate set by the initial-condition measure looks
like.  In the one populated family the campaign could seal, suppression
of time reversal is organized by a Gaussian tail in a measured
effective gap and the Schwinger-shaped variable fails to compete, the
first predicting held-out cells at 0.38 of its own training error and
the second at 3.14 times the first.  The family built to escape that
verdict turned out to contain no events at all, empty at every one of
twelve probed cells, because its alternating slabs cancel between 95
and 99 percent of every deposit with no fixed sign.  Both mandatory
gates then passed their bars on grids that contained none of the events
their claims were about, and both had to be rerun on probe-placed
cells, where conservation holds with every census summing exactly to
twenty thousand and the orientation charge rule failing zero times.
The model carries no quantum phase structure.  A two-level mode swept
twice through an avoided crossing interferes at a level no noise floor
approaches, while the matched classical ensemble carries only the
smooth trend its own pulse overlap predicts, its residual sitting in
the middle of its null.  Reaching that statement took four corrections,
each to an error of the declaration rather than of the model.  A
same-species decay of one particle into two is excluded by parity
rather than missed, since across 441 generic levels on four folding
worldlines the observed branch count changed by minus two, zero, or
plus two and never by anything else.  Declaring one branch class
invisible produced 63 apparent decays with missing energy whose counts
the consumer's own detector model predicted exactly, so apparent decay
here measures what a consumer cannot see.  Lifetimes in this family are
not memoryless, their coefficient of variation measuring 0.159 against
the value one an exponential requires.  The founding question is then
answered under seal for one observable.  A deterministic crossing of a
declared field profile is exponentially suppressed with an exponent
fixed by the distance to the field's nearest complex singularity, four
fitted values agreeing within 1.07 percent across widths spanning a
factor of 2.56, with no distribution anywhere whose tail could be
responsible.  The observable is a per-crossing transfer and not an
event rate, and nothing here is evidence that physical pair creation or
decay works this way.
\end{abstract}

\maketitle

\section{Introduction}
\label{sec:intro}

Take a curve in some space of hidden states and look at it only
through a map into spacetime.  A detector picks a time function on
spacetime, and the observed history of the curve is the sequence of
places where it meets each level of that time function.  If the time
function increases along the curve everywhere, the observer sees one
particle moving forward.  If it has a nondegenerate critical point, the
observer sees one branch turn into three, or none into two, at a single
instant, and the two new branches run in opposite directions of the
detector's time.  Under the standard identification of the direction of
proper time with the sign of charge, that is an apparent pair event
\cite{Stueckelberg1941, Feynman1949}.

This is where the campaign that produced this paper started, and it
separates three claims that must never be run together.  The first is
kinematic.  A nondegenerate fold of the observed time along a smooth
worldline produces exactly two extra branches of opposite orientation,
and the signed count is unchanged.  This is a theorem about smooth maps
and it holds wherever its smoothness hypotheses hold
\cite{Whitney1955, ArnoldGZ}.  The second is dynamical and holds only
inside a declared model.  It asserts that some specified local hidden
dynamics produces such folds at a stable rate that is not put in by
hand.  The third is physical.  It asserts that the rate and the
observables reproduce quantum field theory, including the exponential
suppression Sauter and Schwinger computed for a constant field
\cite{Sauter1931, HeisenbergEuler1936, Schwinger1951, Dunne2004}.  A
successful demonstration of the first is not evidence for the second,
and a successful demonstration of the second in one declared family is
not evidence for the third.  Every section below says which of the
three it is about.

An anti-circularity contract was written before any run and governs all
of them.  No candidate model may contain a command to fold at a wanted
rate, may take the Schwinger exponent as an input probability, may be
fitted to a target rate on the range it is then evaluated on, may
postselect the trajectories that fold or escape or resemble particles,
may use observer-specific coordinates that are not transformed along
with the observer, or may assign energy and charge after inspecting
which branch pairing would be convenient.  What a model may contain is
local fields, a metric or an action, a projection, and a
source distribution fixed before the claim-bearing sweep.  Claims are
carried only by sealed preregistrations whose bars, grids, estimators,
and stopping rules are fixed in advance and whose file hashes are
recorded in a seal ledger before the governed runner executes, which is
the discipline the whole campaign runs under
\cite{ZenodoGames, ZenodoGen2, ZenodoCrypto}.

The arc has one further discipline that it did not start with and had
to earn.  Five separate runs in this arc passed or reported on manifests
that contained none of the events their claims were about, and each
time the error was in the manifest rather than in the instrument.  The
rule that came out of them is now unconditional in this campaign.  Every
registration must cite a committed probe verifying that its manifest
contains the variation its claim is about, for every bound cell and
every bound member, and any arm of any experiment must do the same
even when the arm is a new construction rather than an inherited grid.
Section~\ref{sec:spine} reports all five as record facts and names the
design error in each, together with four bars that failed for reasons
belonging to the declaration rather than to the model and one committed
record that a defective runner produced.

The scope of the whole paper is fixed here and repeated in the ledger
of Sec.~\ref{sec:claims}.  Every dynamical statement below is measured
in a declared finite model with declared constants, integrated
numerically at a declared step, and it is labeled as such.  Two prior
papers of this campaign carry parts of this arc in full detail and are
cited rather than restated, one on the replication of a published
classical pair event \cite{ZenodoCayley} and one on the sealed
Schwinger-scaling negative \cite{ZenodoSchwinger}.  No result reported
here is evidence that physical pair creation or physical particle decay
works this way, no claim about quantum field theory is made or implied,
and the one sealed positive result of the arc is about a per-crossing
energy transfer in a declared one-dimensional model and not about an
event rate.

\section{Simple folds survive perturbation and a bounded model makes
them at a rigid rate}
\label{sec:net}

Let the hidden state be $z(\tau)$ on a smooth manifold, let
$\pi$ carry it to spacetime, and let a detector declare the time
functional $T_u(x) = u_\mu x^\mu$.  An apparent pair event is a
nondegenerate critical point of the observed time,
\begin{equation}
  \frac{d}{d\tau}T_u\bigl(x(\tau)\bigr) = 0,
  \qquad
  \frac{d^2}{d\tau^2}T_u\bigl(x(\tau)\bigr) \neq 0,
  \label{eq:fold}
\end{equation}
near which the observed time has the normal form
$T_u - t_0 = \sigma (\tau - \tau_0)^2$ with $\sigma = \pm 1$.  Crossing
such a point, the unsigned number of preimages of a constant-time slice
changes by exactly two while the signed number, counting each preimage
with the sign of $dT_u/d\tau$, does not change at all.  That is the
kinematic claim, and it is the only statement in this paper with the
label \texttt{[proved]}.

The instrument that finds these points was validated against exact
controls before any dynamics ran, and the tolerances were then frozen
and sealed.  The controls are a monotone null, an exact creation fold,
an exact annihilation fold, a degenerate cubic critical point, and a
double fold, in that order the cases the classifier must find, must
find with the opposite orientation, must refuse, and must find twice
\cite{Whitney1955}.  On the monotone null the instrument reports one
branch and zero critical points at every slice.  On the creation fold
the branch count runs 0, 1, 2 while the signed count stays 0 at every
slice, and on the annihilation fold it runs the other way.  The cubic
control produces one critical point which is classified degenerate and
is never counted as a fold.  The double fold $t = \tau^3 - \tau$
produces band counts 1, 3, 3, 3, 1 with signed count $+1$ everywhere
and both critical points located at $\pm 1/\sqrt{3}$ with the correct
types.  The fitted separation exponent of the square-root law agrees
with $1/2$ to better than $10^{-12}$, and the two independent
implementations agree on branch locations to $4.5 \times 10^{-16}$.

The classifier constants are derived from those margins rather than
chosen.  The first-derivative tolerance sits at $10^{-9}$ against a
worst measured residual of $4.6 \times 10^{-13}$ at located events, and
the second-derivative tolerance sits at $10^{-8}$ against the smallest
genuine fold curvature of $0.208$ found anywhere in the committed
evidence, so no regular point can be read as critical and no fold and
degenerate point can be exchanged.  The freeze states its own
invalidation condition, that a claim-bearing model whose fold
curvatures approach $10^{-5}$ in campaign units forces a re-freeze
before use.

Structural stability was measured next, because a fold that does not
survive perturbation is a coincidence rather than a phenomenon.  The
exact fold was perturbed by bounded random smooth terms at twelve
perturbation norms from $10^{-5}$ to $0.316$, with 500 independent
trials at each.  Fold persistence is exactly 1 at every norm up to
$0.123$ and falls to $0.956$ at the largest, where 22 of 500 trials
lose the fold, and no trial at any norm produced a degenerate
classification.  The median fitted separation exponent stays at
$0.5$ across the whole sweep, departing from it by at most
$2.8 \times 10^{-6}$ at the largest perturbation.  Simple folds persist,
which is what singularity theory predicts.

The negative control measured how a fold rate behaves when nothing
suppresses it.  A bounded local Hamiltonian with an oscillator in the
observed time, an
anharmonic hidden coordinate, and a bilinear coupling produces folds
readily, and the folds are rigid.  A single member integrated by
velocity Verlet at step $10^{-3}$ produces exactly 12 folds, with 7
positively oriented and 6 negatively oriented segments and a relative
energy drift of $1.3 \times 10^{-7}$.  A thousand-member ensemble with
perturbed initial conditions produces exactly 12 folds in every single
member, 12000 folds in total with zero dispersion, at a mean relative
drift of $6.7 \times 10^{-10}$.  The same scenario run on a second
machine in a container reproduces the outcome hash bit for bit.

That uniformity has a cause.  A rate is a property of the
initial-condition measure and the dynamics together, and in a bounded
Hamiltonian with a smooth measure there is nothing to make it
exponentially small.  Exponentially suppressed pair production
therefore requires large-deviation structure somewhere, and the
anti-circularity contract forbids inserting it.  From this measurement
onward the campaign expected a disciplined negative, and said so in
writing before the first sealed run.

\section{A published classical pair event, replicated}
\label{sec:pf3}

Before extending anything the campaign reproduced an existing result.
Stueckelberg-Horwitz-Piron electrodynamics evolves events under a
parameter-dependent five-potential, so the mass shell is not fixed and
$dt/d\tau$ may change sign along one worldline
\cite{Stueckelberg1941, HorwitzPiron1973}.  Land's impulsive Coulomb
scattering configuration was transcribed equation by equation with its
conventions recorded, and its closed-form final velocity agrees with a
numerical solution of the linear system it came from to
$5.7 \times 10^{-14}$, its two published forms of the time component
agree to $1.4 \times 10^{-14}$, and its Rutherford limit reproduces the
published half-angle relation exactly \cite{Land2016}.  The reversal
threshold sits exactly at dimensionless coupling 2 as published.

The replication also produced a finding the campaign kept.  With the
source's own frozen-point convention the kick generator exponentiates
to hyperbolic evolution that never reverses the time velocity, while
the source's impulsive midpoint closure applies the same generator
through a Cayley transform whose pole sits exactly at coupling 2.  The
published threshold is a property of the impulsive closure of the
interaction rather than of continuous accumulation of the same force.
Twenty-seven cells of the sweep have spacelike outgoing velocity, which
is Stueckelberg's own requirement that a reversing worldline cross the
spacelike region twice rather than a transcription error, and the
committed record lists all of them.  Below threshold the smoothed
worldline has zero crossings of the observed time and above it exactly
one, classified an annihilation fold by the sealed classifier, which is
what ties the published pair event to the kinematics of
Sec.~\ref{sec:net}.  The full analysis is a separate paper of this
campaign \cite{ZenodoCayley} and its consequence for what follows is
one clause.  Any candidate whose reversals survive only because its
interactions are closed impulsively is reproducing a Cayley pole and
not a dynamical rate, so every claim-bearing run in this arc integrates
continuously through every interaction.

\section{The Schwinger challenge and its disciplined negative}
\label{sec:pf4}

The first family capable of carrying a physical claim is a
Stueckelberg-type system in which a hidden anharmonic coordinate is
tilted by a Sauter slab,
\begin{equation}
  H = \tfrac12 p_t^2 + \tfrac12 p_u^2
    + \tfrac12 \omega_u^2 u^2 + \tfrac{\lambda}{4} u^4
    + g E L \tanh(t/L)\, u ,
  \label{eq:sauter}
\end{equation}
with $\omega_u = 1.2$, $\lambda = 0.1$, $g = 0.25$, $L = 3$, thermal
initialization in the shifted well at temperature 1, entry at
$t = -4L$, and velocity Verlet at step $10^{-3}$.  A member reverses
when $p_t$ reaches zero, which is the fold of Eq.~\eqref{eq:fold} in
this model.  The question is whether the fraction that reverses is
organized by the Schwinger-shaped combination or by something else.

The pilot answered a prior question first.  Across eighteen cells at
$5 \times 10^4$ members each, with relative drift below
$7 \times 10^{-7}$ in every cell, the family has a critical field above which the
deterministic crossing kick alone reverses the event, and below it
reversal is fluctuation driven.  The transition is steep.  At gap
$0.5$ the reversing fraction moves from $2.8 \times 10^{-4}$ at field
$0.65$ to $0.9115$ at field $0.80$, three decades across a fifteen
percent step in the field, and only two of the eighteen cells landed in
the measurable band at all.  Fixed rectangular grids therefore buy
bounds rather than measurements, and every later manifest in this
family is placed by bisecting a deterministic probe to a declared
target.

Two variables were then declared and compared.  Writing $d(P, E)$ for
the deterministic minimum of the time momentum measured by a
thermal-free probe, which is the effective gap the fluctuation must
cover, and $f$ for the reversing fraction, the two competing models are
\begin{equation}
  -\log f = \alpha + \beta \Bigl(\frac{d}{E}\Bigr)^{2}
  \quad\text{and}\quad
  -\log f = \alpha' + \beta' \frac{d^{2}}{E},
  \label{eq:axes}
\end{equation}
the first a Gaussian tail in the effective gap and the second the shape
the Schwinger exponent has.  Both are fitted by weighted least squares
on training gaps and asked to predict held-out gaps that are disjoint
in the gap variable.

\begin{figure}[t]
\centering
\begin{tikzpicture}
\begin{axis}[paperaxis, height=0.62\columnwidth,
  xlabel={effective-gap variable, $(d/E)^2$},
  ylabel={$-\log f$},
  xmin=0.04, xmax=0.235, ymin=-0.5, ymax=10.8,
  xtick={0.05,0.10,0.15,0.20}, scaled x ticks=false,
  x tick label style={/pgf/number format/fixed,
    /pgf/number format/precision=2},
  legend pos=north west, legend cell align=left]
\addplot[gray!70!black, densely dotted, no marks, domain=0.04:0.235,
  samples=2] {-0.12429061298285399 + 43.176187880553684*x};
\addlegendentry{effective-gap fit, training cells}
\addplot[only marks, mark=*, mark size=2.1pt, blue!70!black]
  table[x=xg, y=nlogf] {figdata/pf4_train.dat};
\addlegendentry{training gaps 0.55, 0.75, 0.95}
\addplot[only marks, mark=square*, mark size=2.1pt, red!70!black]
  table[x=xg, y=nlogf] {figdata/pf4_held.dat};
\addlegendentry{held-out gaps 0.65, 0.85}
\end{axis}
\end{tikzpicture}
\caption{The sealed axis comparison of PREREG-PF4-002.  Each point is
one probe-placed cell of $5 \times 10^4$ members, plotted at the
declared effective-gap variable against the measured suppression.  The
dotted line is the weighted fit made on the training cells alone.  The
held-out cells lie on it, and the model's held-out weighted error is
0.38 of its own training error against a declared bar of 2, while the
Schwinger-shaped variable of Eq.~\eqref{eq:axes} is 3.14 times worse on
the same held-out cells against the same bar.  Ten of twelve training
and seven of eight held-out cells met the sealed usability rule and are
shown.  All values are read from the committed evidence record.}
\label{fig:pf4}
\end{figure}

The first sealed attempt returned the third branch its own
preregistration had defined.  All twenty probe-placed cells landed in
the measurable band, both prescription controls agreed per member
exactly at $z = 0.00$ under fourth-order Runge-Kutta and under step
halving, the zero-field null produced zero reversals, and the fitted
effective-gap slope of $41.24$ confirmed the pilot's two-point
prediction.  The effective-gap model beat the Schwinger-shaped model by
a factor of $3.28$ on held-out weighted error.  But the declared bar
required a held-out weighted mean squared error of at most 4 and the
measurement was $81.22$, because at that ensemble size the residual gap
dependence of the suppression at fixed effective-gap variable is many
binomial standard errors.  The verdict was neither-axis, and the
recorded reading is that the bars had been written in binomial units
and were therefore testing a model nobody had claimed.

The second attempt declared in writing, before it existed, that it
would set every bar against the model's own training inadequacy rather
than against binomial noise, and it passed both.  The effective-gap
model's held-out weighted error is $0.3842$ of its training error
against a bar of 2, so it generalizes to gaps it never saw better than
it fits the gaps it trained on, and the Schwinger-shaped model is
$3.1398$ times worse on the same held-out cells against the same bar.
The fitted effective-gap slope is $43.18$, ten of twelve training and
seven of eight held-out cells were usable under the sealed rule, both
prescription controls again agreed at $z = 0.00$, and the zero-field
null again produced zero reversals.  Figure~\ref{fig:pf4} shows the
comparison.

The label this earns is \texttt{[demonstrated-in-model]} and its
content is narrow.  In this declared family the suppression of time
reversal is a Gaussian tail of the initial-condition measure in a
measured effective gap, and the Schwinger-shaped variable fails to
compete.  That is a negative about a model family.  It is not a
statement about other families and it is not a statement about physical
pair production.  What it left open is the construction question, which
became the summit of the whole arc.  Does any constructible family have
an exponent that is not a measure tail?  The detailed treatment of this
negative is a separate paper \cite{ZenodoSchwinger}.

\section{The family built to answer the summit was empty}
\label{sec:pf4b}

A measure tail exists because there is a measure.  The obvious way to
remove it is to remove the ensemble, so the next family is a single
deterministic event crossing an alternating train of Sauter slabs,
\begin{equation}
  \mathrm{tilt}(t) = g E L \sum_k (-1)^k
    \tanh\!\Bigl(\frac{t - kD}{L}\Bigr),
  \label{eq:train}
\end{equation}
at field $0.6$, slab length 3, spacing 18, integrated continuously by
fourth-order Runge-Kutta.  Each crossing deposits an
analyticity-suppressed residual, the deposits accumulate, and the
observable is the slab count at reversal.  With one declared initial
condition there is no distribution anywhere, so an exponential in this
family could not be a measure tail.

The sealed run sealed itself unevaluable.  Every one of the four
training and three held-out velocities capped at 4000 slabs without a
single reversal, and so did every cell of the mandatory planted
control, so no fit existed to read and the record carries a
manifest-design failure rather than a pass or a refutation.  The
diagnosis was that the manifest had been placed from a probe of a
different observable at different constants, and the repair was a
direct-observable probe run with the sealed dynamics verbatim.  That
probe measured the reversal count at twelve cells, three field
strengths crossed with four velocities, and found zero reversals in all
twelve at a cap of 300 slabs.  The template registration that would
have succeeded PF4-003 therefore cannot bind, because its own
verification clause requires every manifest cell to have reversed below
half the probe cap, and no cell reversed at all.

A separate study measured why, and its declared compound flag read
False in a way that separates cleanly.  The declared hypothesis had two
halves, that the alternating slabs cancel the deposits and that the
extrapolated slab count to reversal exceeds the cap.  Cancellation is
confirmed and it is strong.  The net drift per slab is $3.3$, $0.9$,
and $5.0$ percent of the mean deposit magnitude in the three declared
cells, so between 95 and 99 percent of every deposit is undone by the
slabs that follow, with the sign of consecutive deposits flipping on
$54$, $68$, and $27$ percent of steps.  The second half failed because
it was written against the absolute drift, and the third cell's drift
is positive at $+0.0096$ per slab, carrying the trajectory away from
reversal rather than toward it.  The corrected reading is stronger than
the declaration.  The pulse train does not merely deposit too slowly.
Its residual after cancellation has no fixed sign, and at the largest
declared field it accelerates the event away from the fold, which is
why every probed cell is empty.

\section{Two mandatory gates, each of which had to be run twice}
\label{sec:gates}

The campaign sequence makes conservation and covariance mandatory gates
before any further claim, not optional follow-up analyses.  Both gates
were built, validated on exact controls, sealed, run, and passed.  Both
first runs then turned out to have certified nothing about the events
they were about, and both had to be redeclared on probe-placed
manifests.

The conservation instrument provides a complete census in which every
trajectory ends in exactly one declared class, numerically failed,
reversing, transmitted, or capped, so the missing-trajectory count is
zero by construction and every reported statistic carries an explicit
denominator with a flag that fires whenever a deletion would change a
number.  It tracks the conserved Hamiltonian every twenty-five steps
for every member including failed members up to their failure step.
And it declares the orientation charge to be the sign of $dt/d\tau$
before any claim run, with the conservation statement written as a
path-degree rule, that the signed crossing count of any generic level
equals the endpoint bookkeeping, which is exactly the statement that
folds make orientation charge only in cancelling pairs.  Its controls
all pass.  A harmonic control drifts by $4.4 \times 10^{-7}$ against a
control bar of $10^{-6}$, a field-free cell classifies 2000 of 2000
transmitted, a deterministically reversing cell classifies 2000 of 2000
reversing, an absurd step size counts 200 of 200 numerically failed
with zero missing, a short cap counts 100 of 100 capped, the analytic
double fold obeys the path-degree rule at all 41 generic levels, and
the deliberate deletion trap is mechanically flagged.

The first sealed gate run passed every bar and tested nothing.  Its
eight cells are all evaluable, every census sums exactly to 20000 with
zero missing trajectories, the largest relative energy residual across
cells is $6.5 \times 10^{-7}$ against the sealed bar of $10^{-5}$, and
the path-degree rule holds at every generic level on all 96 declared
polyline members with zero failures and zero refused levels.  The same
census also records zero reversing trajectories in every cell.  All
160000 members transmitted.  The sealed claim was about the accounting
of apparent pair events and this grid contained none, so what the run
certifies is the accounting machinery on transmitted worldlines.  The
named error is the manifest.  A product grid of gaps against fields was
declared where the sealed family's own cells had been placed by
bisection, which is how the manifest landed at fields where no member
reverses.

The repair binds cells from a committed placement probe that bisects
the field at each declared gap until the reversing fraction lands
inside a declared window, measuring $0.1000$, $0.0970$, $0.0995$, and
$0.1030$ at the four bound cells, and it adds one bar that makes a
vacuous pass impossible, that every evaluable cell must record a
reversing count above zero and the run must record at least two
thousand reversals in total.  The rerun passes all of it.  The four
cells record $1998$, $1952$, $1881$, and $2237$ reversing members out
of 20000 each, 8068 audited apparent pair events in total, every census
exact with zero missing, the largest relative energy residual
$7.3 \times 10^{-7}$ against the same $10^{-5}$ bar, and zero
path-degree failures and zero refused levels on all 48 declared
polyline members.  Conservation and complete accounting hold on cells
that contain the events.

The covariance instrument audits five things, and four of them applied
here.  Translation covariance shifts the field profile and the initial
data together and requires every event worldpoint to move by exactly
the offset, measured at $2.7 \times 10^{-13}$ on its validation
members.  Reparametrization invariance scales the Hamiltonian and
requires the per-worldline event count to be exactly invariant while
the rate per unit parameter scales by the declared factor, which it
does exactly, and the witness this freezes is that the physical rate of
this model class is per worldline and never per parameter tick.  The
boost audit runs on the family with genuine hyperbolic structure, the
replication of Sec.~\ref{sec:pf3}, where flow and boosts commute at
$4.5 \times 10^{-13}$, the invariant is conserved at
$7.1 \times 10^{-11}$, and the future-cone margin stays positive across
the declared rapidity grid, so the never-reverses finding of that
family is boost invariant.  The observer audit measures the event count
twice for a declared family of detector time functionals, once from the
observer's own record and once from the model's detector prediction,
and requires exact agreement at every declared observer.  The gauge
clause is recorded not applicable rather than passed, because the
Sauter tilt is declared non-electromagnetic and has no field tensor,
and the gauge machinery stands validated on a separate electromagnetic
control at $2.2 \times 10^{-13}$ for the physical history and
$1.1 \times 10^{-12}$ for the predicted canonical shift.

The first sealed covariance run returned vacuous, and its own
anti-vacuity bar is what caught it.  The declared deterministic member
at each bound cell produces zero folds, so four cells were audited for
the covariance of events that were not there.  The named error is the
member rather than the cell.  The thermal ensemble reverses about ten
percent of the time at these cells, so the members that reverse carry
transverse momentum in the tail while the declared member did not.  A
committed member probe then scanned the initial transverse momentum at
each cell and bound the members that fold, with verified fold counts 4,
2, 2, and 2.  The rerun passes.  Fold counts are exactly the probe's,
10 folds in total and every count even, translation shifts move the
first-fold worldpoint by exactly the offset to
$3.1 \times 10^{-14}$, the reparametrized rate ratio deviates from the
declared factor by exactly zero, the boost audit repeats its values,
and the observer count read from the observer's own record equals the
detector-model prediction at every one of the six declared observers on
every member.  Observer-dependent particle number survives in this
model, lawfully, because a declared detector model predicts every count
exactly.

\section{The model carries no quantum phase structure}
\label{sec:pf7}

With both gates cleared the ceiling test is ungated.  Its question is
whether the fold model reproduces more than a classical worldline
picture, and its discriminator was declared with the reasoning that
makes it refutable.  A quantum mode swept twice through an avoided
crossing accumulates a relative phase between the two passages, so its
excitation probability oscillates with the delay between them, which is
Landau-Zener-Stueckelberg interferometry
\cite{Landau1932, Zener1932, Stueckelberg1932, Shevchenko2010}.  A
classical ensemble of independent trajectories has no phase to
accumulate.  Both arms are reduced to one statistic, the largest
nonzero-frequency power of the mean-removed discrete Fourier spectrum
of the curve divided by its total power.

The quantum arm is a two-level mode with declared gap $0.30$ and sweep
rate 1, integrated by fourth-order Runge-Kutta, whose single-passage
control must match the closed form
\begin{equation}
  P = \exp\Bigl(-\frac{\pi\,\Delta^{2}}{2 v}\Bigr) = 0.868167
  \label{eq:lz}
\end{equation}
at the declared gap and sweep rate.  The classical arm is the pilot
fold dynamics with the single slab replaced by two Sauter pulses
separated by the same delay, with the reversing fraction of a seeded
thermal ensemble as its observable.

The first run failed two instrument bars and its apparent ceiling
finding was vacuous.  The single-passage control measured $0.8631$
against Eq.~\eqref{eq:lz}, a deviation of $5.06 \times 10^{-3}$ against
a declared $2 \times 10^{-3}$, because the declared sweep range of
forty natural units does not reach the asymptotic limit at that gap.
The quantum oscillation statistic measured $0.2238$ against a declared
$0.30$ while the recorded curve plainly oscillates, because the
Stueckelberg phase between crossings at plus and minus a half-delay is
the integral of the detuning between them and is therefore quadratic in
the half-delay, so a grid uniform in the delay undersamples a chirp and
spreads its power across bins.  The norm of the two-level state is
conserved to $5.5 \times 10^{-14}$, so the mode holds at most one
excitation exactly.  And the classical arm returned a reversing
fraction of exactly $1.0$ at every one of the twenty-four delays, all
2000 members of every cell reversing, so its oscillation statistic was
computed on a constant curve and was zero for a reason with nothing to
do with interference.  The declared field was strong enough that two
pulses reverse everything.

\begin{figure}[t]
\centering
\begin{tikzpicture}
\begin{groupplot}[group style={group size=1 by 3, vertical sep=32pt},
  paperaxis, height=0.36\columnwidth]
\nextgroupplot[xlabel={squared half-delay},
  ylabel={excitation probability},
  xmin=14, xmax=102, ymin=-0.03, ymax=0.46]
\addplot[blue!70!black, mark=*, mark size=1.5pt]
  table[x=t2, y=p] {figdata/pf7_quantum.dat};
\nextgroupplot[xlabel={delay between pulses},
  ylabel={reversing fraction},
  xmin=7.6, xmax=20.4, ymin=0.470, ymax=0.524,
  scaled y ticks=false,
  y tick label style={/pgf/number format/fixed,
    /pgf/number format/precision=3}]
\addplot[red!70!black, mark=square*, mark size=1.5pt]
  table[x=d, y=f] {figdata/pf7_classical.dat};
\nextgroupplot[ybar, bar width=11pt, ylabel={oscillation statistic},
  xtick={1,2,3,4,5},
  xticklabels={null mean, null 95th, classical raw,
    classical residual, quantum},
  x tick label style={rotate=28, anchor=east, font=\footnotesize},
  xmin=0.4, xmax=5.6, ymin=0, ymax=1.02,
  nodes near coords,
  every node near coord/.append style={font=\footnotesize,
    /pgf/number format/fixed, /pgf/number format/precision=4}]
\addplot[fill=blue!35, draw=blue!70!black]
  table[x=idx, y=stat] {figdata/pf7_stats.dat};
\end{groupplot}
\end{tikzpicture}
\caption{The ceiling measurement after its four corrections.  The top
panel is the quantum arm's excitation probability on the grid uniform
in the squared half-delay, where the chirp becomes a single frequency.
The middle panel is the classical arm's reversing fraction at eight
thousand members per delay, whose whole range is 0.0425 and which falls
smoothly as the two pulses stop overlapping.  The bottom panel places
the declared statistic against the null distribution built from the
same binomial counts, showing the null mean, the null ninety-fifth
percentile, the raw classical value, the classical value after the
declared smooth overlap model is removed from the curve and from every
null curve alike, and the quantum value.  The classical residual sits
below the null mean at a p-value of 0.757 and the quantum value sits
past every null draw.  All values are read from committed evidence
records.}
\label{fig:pf7}
\end{figure}

Three corrections followed, each to a named error.  The sweep range
became 120 natural units, the delay grid became uniform in the squared
half-delay, and the classical field was bound from a committed probe
that bisects it to a reversing fraction near one half and then verifies
the fraction at the extreme delays.  That probe rejected the originally
declared delay range by its own window rule, because at the shortest
delay the two pulse centres sit closer together than the pulse length
and the pulses act as one stronger pulse, returning a fraction of
exactly 1.  A second probe measured $0.5245$, $0.4945$, $0.4750$,
$0.4835$, and $0.4955$ across the candidate range, all inside the
declared window, and the corrected run bound that range.

The corrected run passed all three instrument bars.  The single-passage
control now measures within $1.49 \times 10^{-3}$ of Eq.~\eqref{eq:lz},
the quantum oscillation statistic rose from $0.2238$ to $0.8702$ under
the grid uniform in the squared half-delay, which confirms that the
earlier failure was the grid rather than the physics, and the norm is
conserved to $5.9 \times 10^{-14}$.  The classical arm varies as the
probes promised, between $0.4645$ and $0.5245$ across thirty-two
delays.  The ceiling finding then failed its own declared bar.  The
classical statistic measured $0.3626$ against a declared $0.15$ and the
quantum-to-classical ratio $2.40$ against a declared 3.

That bar was unattainable as declared, and the reason belongs to the
statistic.  On a thirty-two point curve the statistic is the largest of
sixteen powers divided by their sum, which is a sizable fraction even
for pure noise, and each cell carries binomial noise of about $0.011$
from two thousand members while the whole curve spans only $0.06$.  The
next run measured that floor instead of asserting it, drawing twenty
thousand null curves from binomial counts at the curve's own mean.  The
null mean is $0.2165$ and its ninety-fifth percentile is $0.3310$, so
the declared bar of $0.15$ sat below the statistic's floor and could
not have been met by any curve of that length.  Against that null the
classical statistic has a p-value of $0.0262$ and the quantum statistic
a p-value of $0.0$.  The protocol had declared in advance that a small
classical p-value would be escalated rather than read as a ceiling, and
it was.

The escalation quadrupled the ensemble to eight thousand members per
delay and tested a declared mechanism, that the hidden oscillator rings
at its own period and meets the second pulse in or out of phase with
that ringing.  The structure survived the larger ensemble at a p-value
of $0.02605$, essentially unchanged from $0.0262$ at a quarter of the
members, so it is not shrinking the way a fluctuation would.  The
mechanism was refuted.  The spectral peak sits at a period of $12.52$
in delay while the declared ringing period is $5.236$, and $12.52$ is
the span of the delay range, so the peak is the lowest nonzero
frequency bin, which is where a monotone trend lands rather than an
oscillation.  The curve falls from a mean of $0.4942$ over the first
half of the range to $0.4829$ over the second.  Two pulses of length
three separated by eight still overlap and by twenty do not, so the
combined kick weakens smoothly with delay.  The consequence recorded
with it is that the declared statistic counts a trend as structure, so
neither p-value was evidence about oscillation.

Removing a declared linear trend from the curve and from every null
curve alike then failed its own first item and refuted the diagnosis
that produced it.  The classical statistic rose from $0.4480$ to
$0.5014$ and its p-value fell from $0.02605$ to $0.01185$, while the
quantum statistic was untouched at $0.8702$ before and after, as
predicted for a peak sitting at a high frequency.  The trend itself is
real, a drop of $0.0234$ across the range against a standard error of
$0.0056$, which is $4.19$ standard errors.  The arithmetic of the rise
is plain, since removing power from the lowest bin shrinks the total
and a peak elsewhere takes a larger share.  What the failure showed is
that a linear model does not absorb an exponential approach to an
asymptote, and that the null, which draws every cell at one common
mean, counts any smooth departure from a straight line as structure.

The last refinement declared its model from the mechanism rather than
fitting it blind.  Two pulses of declared length three overlap by an
amount falling exponentially with their separation, so the declared
model is an asymptote plus an exponential approach at that length,
linear in its two free coefficients, removed from the curve and from
every null curve alike.  It passed all three items and both findings
hold.  The classical statistic drops from $0.4480$ to $0.2122$ with a
p-value of $0.757$, which sits in the middle of the null rather than in
its tail, the quantum statistic is $0.8724$ after the same treatment,
and the trend remains real at $4.19$ standard errors.
Figure~\ref{fig:pf7} shows the corrected picture.

The terminus was declared before that run so that a fourth refinement
could not become a search for a wanted answer, and the reading it fixed
is the one that holds.  The classical arm carries exactly the smooth
delay dependence its own mechanism predicts, and once that mechanism is
accounted for there is nothing left in its spectrum that noise does not
explain.  The quantum arm interferes at a level no noise floor
approaches.  The fold model is capped as a classical representation.
It does not carry the phase structure that produces interference, and
this caps the model without touching the kinematic claim of
Sec.~\ref{sec:net} or any measured result of this arc.

\section{Decay is excluded by parity and imitated by blindness}
\label{sec:pf8}

A fold changes the observed branch count by exactly two.  A
same-species decay of one particle into two is a change of exactly one.
Those two sentences settle the decay question for complete continuous
worldlines under generic projections before any run, and the run
measures whether the model obeys its own arithmetic and what a
restricted observer sees instead.

The design had assumed the pair-creation configuration, in which the
worldline begins and ends on one side of the observed level so the
branch count is even.  A committed probe measured that configuration
absent on every bound member, since all four members end at their own
largest observed time and begin at their own smallest, making both
candidate windows exactly zero wide.  The bound members are
through-going, one branch in and one branch out, which is the
configuration a decay question wants anyway.  On a through-going
worldline the complete observer's unsigned branch count is odd at every
generic level, and a fold changes it by two, so a sweep reads one
branch, then three, then one.  The reframing was declared before the
run.  The same probe fixed the level ladder, because it measured the
fold excursions to be as narrow as $0.000253$ in observed time against
a uniform ladder spacing near $0.2$, so a uniform ladder steps over the
three-branch fibers almost everywhere.  The declared ladder merges a
uniform background with nine levels inside each interval spanned by a
consecutive pair of fold times.

\begin{figure}[t]
\centering
\begin{tikzpicture}
\begin{axis}[paperaxis, ybar, bar width=9pt, height=0.56\columnwidth,
  xlabel={branches recorded at a generic level},
  ylabel={observations},
  xtick={0,1,2,3}, xmin=-0.6, xmax=3.6, ymin=0, ymax=540,
  legend pos=north east, legend cell align=left,
  nodes near coords, every node near coord/.append style={
    font=\scriptsize, /pgf/number format/1000 sep={}}]
\addplot[fill=blue!35, draw=blue!70!black]
  table[x=n, y=count] {figdata/pf8_complete.dat};
\addlegendentry{complete observer}
\addplot[fill=red!30, draw=red!70!black]
  table[x=n, y=count] {figdata/pf8_blind.dat};
\addlegendentry{declared blind consumer}
\end{axis}
\end{tikzpicture}
\caption{What two observers of the same four worldlines record, pooled
over all 441 generic levels of the declared ladder.  The complete
observer records only 1 or 3 branches, which is the parity fixed by the
worldline endpoints, and every change between adjacent levels is minus
two, zero, or plus two.  The consumer declared blind to branches whose
transverse coordinate lies in a fixed window records 0, 1, 2, or 3, so
its parity is not conserved and 63 of its transitions are a rise by
exactly one, which is what a decay of one particle into two looks like.
Every one of its 441 counts was predicted exactly by its own detector
model.  All values are read from the committed evidence record.}
\label{fig:pf8}
\end{figure}

The complete-observer arm passed with zero odd events.  Four
deterministic members were audited on 441 generic levels and 400 kept
thermal polylines on five levels each.  The complete observer's
unsigned branch count took only the values 1 and 3 and every change
between adjacent levels was minus two, zero, or plus two, giving zero
odd events and zero parity failures on either arm, with fold counts 4,
2, 2, 2, all even.  The frozen path-degree rule reported zero failures
on all 441 deterministic levels and all 2000 thermal polyline checks,
and the signed count on every member was exactly $+1$ at every level
with spread zero, so no fold moved it.  Folds make orientation charge
only in cancelling pairs, measured rather than assumed.

The blind consumer is where the phenomenology appears.  The declared
consumer fails to record any branch whose transverse coordinate at the
crossing lies in a closed window fixed from the probe record before the
run, which the probe measured to hide 147 of 533 pooled crossings,
$27.6$ percent, some but not all.  Across 441 observations that
consumer failed to record 147 crossings and read a parity different
from the complete observer on 129 of them, a flip fraction of
$0.2925$.  Sixty-three of those were apparent decay events, meaning
adjacent levels across which the blind count rose by exactly one, one
particle becoming two with missing energy.  Every single apparent count
matched the detector-model prediction exactly, zero mismatches on 441
observations, computed by an independent route from the consumer's own
record.  On the member where all three branches of the zigzag lie
inside the window the consumer sees nothing at all where the complete
observer sees three.  Figure~\ref{fig:pf8} shows both observers.
The anti-vacuity clause holds throughout, every member folding, the
ladder resolving 46 multibranch levels, the consumer hiding 147
crossings, and the thermal census containing 1963 reversing members
with the census summing exactly to 20000 and a largest relative energy
residual of $6.3 \times 10^{-7}$.

Lifetimes were measured with an expected negative declared in advance.
Real decay is memoryless, so the survival curve of a memoryless process
has a coefficient of variation of exactly one.  Re-running the frozen
census at a ladder of step caps with the seed held fixed gives the
exact cumulative distribution of the first-reversal time over the same
members.  Pooling the four cells, 1963 first-reversal times have mean
proper time $18.92$ and standard deviation $3.01$, a coefficient of
variation of $0.159$.  Per cell it is $0.0526$, $0.0530$, $0.0694$, and
$0.0764$, and the pooled figure is larger only because the four cells
have different means.  The largest deviation between the empirical
survival curve and the exponential whose mean matches it is $0.506$
pooled and near $0.57$ in every cell, most of the available range, and
shifting to the measured onset does not rescue it, giving a shifted
coefficient of variation of $0.494$ and a shifted deviation of
$0.232$.  An exact cross-check on 32 first-reversal steps read from
kept polylines gives mean $18.94$ and coefficient of variation
$0.188$.  First-reversal time in this family has a hard onset set by
when the trajectory reaches the slab and a narrow spread after it,
which is not a hazard at all, and it is the negative control's lesson
again, that fold statistics are borne by the declared measure.

The answer the arc records is therefore two-sided.  For the decay
everyone means, no.  A complete observer of a continuous worldline
under a generic projection cannot see one particle become two, and that
is parity rather than a limit of the search.  For a decay nobody can
distinguish from it without seeing every branch, yes.  The parity-even
zigzag of a parent plus a created pair reads as a decay with missing
energy the moment one branch class goes unrecorded, at counts the
consumer's own detector model predicts exactly.  Apparent decay in this
model is a measurement of consumer blindness and never a property of
the worldline.  This is the same structure the campaign's entropy track
measured as an invisible leak \cite{ZenodoEntropy}, and it is the
phenomenology of decays with unobserved products without any of the
physics.  The caveats belong beside it.  This is four members and one
declared blind window, so the counts of apparent decays are counts on
one ladder and not a rate, the kept thermal polylines are half
worldlines because the frozen census stops a member at its first
reversal, and nothing in this section is sealed.

\section{The summit question, answered under seal for one observable}
\label{sec:summit}

The open question left by the sealed negative was whether any
constructible family has an exponent that is not a measure tail.  A
single deterministic crossing has one declared initial condition, so
there is no distribution whose tail could produce a suppression, and an
exploratory probe had already measured the per-crossing energy transfer
of such a crossing to be exponential in the inverse velocity with a
coefficient of determination of $0.9921$, at a slope of $-7.8186$
against a naive estimate of $-11.3097$.  What that probe could not say
is whether the exponent is dynamical or an accident of a fitted range.

The decisive property is analyticity.  An exponent set by the field's
nearest complex singularity is a dynamical quantity, because no measure
tail knows where a function's poles are, and the adiabatic
exponentiation of a two-level crossing is controlled exactly by that
distance \cite{Dykhne1962, DavisPechukas1976}.  Writing $d$ for the
distance from the real axis to the nearest complex singularity of the
force profile and $P$ for the crossing velocity, the natural variable
is the adiabaticity
\begin{equation}
  \kappa = \frac{\omega_u\, d}{P},
  \label{eq:kappa}
\end{equation}
and in that variable the parameter-free estimate for the log transfer
is a slope of exactly $-2$.  The observable throughout is the residual
oscillator energy about the instantaneous well at exit, which is well
defined for every profile tested including the ones with algebraic
tails.

The first run declared its cells at fixed velocities and failed three
of four bars.  The timestep control passed at
$6.9 \times 10^{-11}$, so nothing in it was numerical, and the
sech-squared cell at width three reproduced the earlier probe's slope,
$-7.8154$ against $-7.8186$, an independent replication by a separate
runner.  The Lorentzian profile behaved as an analyticity-controlled
exponent should, with coefficients of determination $0.99997$,
$0.99991$, and $0.99933$ and measured exponents $1.0095$, $1.0344$, and
$1.0658$ times the naive estimate, while the sech-squared profile gave
$0.5430$, $0.6910$, and $0.8098$ and drifted with width.  The named
error is the grid.  Velocities were declared in absolute terms and
applied to every width, so the three widths were compared at different
adiabaticities and the narrowest cells sat partly outside the range of
adiabaticity where the exponential form holds.

Placing each cell at a declared adiabaticity instead removed the drift
entirely and refuted a declared expectation.  The Lorentzian exponents
in the variable of Eq.~\eqref{eq:kappa} are $-2.0661$, $-2.0639$, and
$-2.0607$ across a factor of two in width, a spread of $0.26$ percent,
each sitting $3.3$ percent above the parameter-free value, at
coefficients of determination of $0.9998$.  The sech-squared exponents
are $-1.0490$, $-1.0478$, and $-1.0462$, a spread of $0.27$ percent.
The declared expectation that the sech-squared profile is not
width-universal is refuted, and the earlier 39 percent spread was
entirely an artifact of the grid.  Both profiles carry an exponent
fixed by the adiabaticity built from their own singularity distance.
What they do not share is the constant, the Lorentzian sitting at the
naive value and the sech-squared at $0.5245$ of it.  The run's verdict
is nonetheless FAIL, on a bar requiring a coefficient of determination
of $0.99$ applied to both profiles when only one of them was predicted
to be exactly exponential, and the sech-squared cells sit at $0.9886$,
$0.9887$, and $0.9889$.

Applying the gate content to the new profile then found a defect in the
observable's definition before anything about it was sealed.  The
conserved quantity of the extended
system drifts by at most $8.6 \times 10^{-14}$, the census is complete
with every crossing transmitted and no nonfinite member, and
translation covariance is exact to $8.8 \times 10^{-14}$, so the
transfer does not care where the field is centred.  But moving the
entry point from twenty widths to thirty changes the transfer by
$1.289$, $1.295$, and $1.291$ percent at the three widths against a bar
of $10^{-8}$.  The diagnosis is that the arctangent tilt approaches
its asymptote algebraically rather than exponentially, so the field
never fully turns off and a transfer measured from any finite entry
carries a tail contribution.  A second clause missed for a different
reason, halving the timestep changing the transfer by about
$3 \times 10^{-7}$ against the same $10^{-8}$, which is a bar set
tighter than the declared timestep can deliver and is a fault in the
declaration.

The span was then measured rather than assumed.  Along a declared
ladder of entry spans the Lorentzian transfer's successive relative
changes are $1.278$, $0.586$, $0.130$, and $0.015$ percent, falling
steadily, while the sech-squared transfer changes by at most
$1.1 \times 10^{-13}$ along the same ladder against a declared bar of
$10^{-9}$, so the effect belongs to algebraic tails and not to the
method.  The exponent is far more robust than the transfer, changing by
at most $0.23$ percent between the nearest and furthest spans, and the
width universality survives at the furthest span with a spread of
$0.36$ percent.  The tail contribution lands in the prefactor and
leaves the exponent alone.  Those two exponent figures are read from
the corrected record of this run.  The record first committed for it
was written by a defective runner and is superseded, which
Sec.~\ref{sec:spine} reports with the rest of the failure record.

\begin{figure}[t]
\centering
\begin{tikzpicture}
\begin{axis}[paperaxis, height=0.62\columnwidth,
  xlabel={declared profile width},
  ylabel={fitted exponent in the adiabaticity},
  xmin=1.25, xmax=4.95, ymin=-2.20, ymax=-0.95,
  legend style={font=\scriptsize, draw=gray!45, fill=white,
    at={(0.5,1.02)}, anchor=south, legend columns=2,
    /tikz/every even column/.append style={column sep=5pt}},
  legend cell align=left]
\addplot[gray!70!black, densely dotted, no marks, domain=1.25:4.95,
  samples=2] {-1.9200522391689598};
\addlegendentry{sealed mean}
\addplot[only marks, mark=*, mark size=2.4pt, blue!70!black]
  table[x=L, y=slope] {figdata/pf4_009_sealed.dat};
\addlegendentry{sealed, simple pole}
\addplot[only marks, mark=triangle*, mark size=2.6pt, teal!70!black]
  table[x=L, y=slope] {figdata/pf4_008c_gud.dat};
\addlegendentry{exploratory, simple pole}
\addplot[only marks, mark=diamond*, mark size=2.6pt, orange!80!black]
  table[x=L, y=slope] {figdata/pf4_005b_lorentz.dat};
\addlegendentry{exploratory, branch point}
\addplot[only marks, mark=square*, mark size=2.0pt, red!70!black]
  table[x=L, y=slope] {figdata/pf4_008c_sech2.dat};
\addlegendentry{exploratory, double pole}
\end{axis}
\end{tikzpicture}
\caption{Every fitted exponent of the summit sequence, plotted against
the declared width of its field profile.  The four sealed values span
widths differing by a factor of 2.56 and agree within 1.07 percent of
their mean, which the dotted line marks, against a declared bar of two
percent, and that agreement is the sealed
claim that the exponent depends only on the adiabaticity built from the
singularity distance and not on the width.  The two exploratory
families whose nearest singularity is first order sit near the same
value while the family whose nearest singularity is a double pole at
the same distance sits at about 1.1, so distance alone does not fix the
exponent.  All values are read from committed evidence records.}
\label{fig:summit}
\end{figure}

Separating the singularity's order from its tail behaviour was the last
exploratory step, and it was declared with two mutually exclusive
hypotheses so the run would discriminate rather than confirm.  A
Gudermannian tilt has a first-order pole at the same distance as the
sech-squared family's double pole and has exponentially decaying tails,
so the comparison holds the distance fixed and varies the order.  The
entry-point clause the Lorentzian profile failed at
$1.3 \times 10^{-2}$ is passed by this one at $5.7 \times 10^{-11}$,
with the timestep and translation clauses at $4.9 \times 10^{-13}$ and
$6.5 \times 10^{-15}$.  The discriminator answered for the order.  The
Gudermannian exponents are $-1.9333$, $-1.9286$, and $-1.9217$, a
spread of $0.61$ percent, while the sech-squared exponents on the
identical grid are $-1.1028$, $-1.1015$, and $-1.0997$, a spread of
$0.29$ percent, a ratio of $1.75$.  Two profiles whose nearest
singularities sit at the same distance and differ only in order carry
different exponents, so the distance alone does not fix it.

That run nonetheless failed on conservation, drifting at
$1.9 \times 10^{-10}$ against a bar of $10^{-10}$, and the declared
repair refuted itself.  Halving the timestep left the drift at
$2.0 \times 10^{-10}$, so the drift was not integration error.  The
diagnosis is a floating-point inconsistency in the tilt.  A hyperbolic
tangent of twenty rounds to exactly one in double precision while its
own derivative is still nonzero at that point, so the tilt and the
force become mutually inconsistent in the tails, and because the
Hamiltonian carries the tilt the inconsistency reads as an apparent
conservation violation.  The sech-squared profile is immune because its
force at the same point is far smaller.  Evaluating the tilt so that it
saturates together with its derivative dropped the worst drift to
$2.6 \times 10^{-12}$ and moved the exponents only in their seventh
digit, exactly as the diagnosis predicted, and that run passes all
eight of its items with entry-point independence at
$3.0 \times 10^{-9}$, timestep independence at $6.6 \times 10^{-13}$,
and translation covariance at $2.5 \times 10^{-15}$.

The sealed registration then bound the Gudermannian profile on a
manifest disjoint from every prior run, four widths spanning a factor
of $2.56$ crossed with six declared adiabaticities, with every bar set
at the precision the prior measurements support rather than at the
precision they appeared to have.  It passes all six.  The four fitted
exponents are $-1.929735$, $-1.924383$, $-1.916823$, and $-1.909268$, a
spread of $1.07$ percent against a declared bar of two percent.
Coefficients of determination are $0.9989$ or better in every cell.
The extended Hamiltonian drifts by at most $3.7 \times 10^{-13}$
against a bar of $10^{-11}$, entry-point independence measures
$9.2 \times 10^{-9}$ against $10^{-7}$, every transfer lies inside the
declared signal window, and no cell is nonfinite and none folds.
Figure~\ref{fig:summit} places those four values beside every
exploratory exponent of the sequence.

The summit question has a sealed answer for one observable.  A
constructible family has a suppression whose exponent depends only on
the adiabaticity built from the distance to the field's nearest complex
singularity, and is therefore independent of the width.  Each crossing
is deterministic with one declared initial condition, so no
distribution exists whose tail could produce the suppression.  The
exponent is dynamical and it is not a measure tail.

The registration's own non-claims bind that reading and are repeated
here because they are the whole scope of the result.  The observable is
the per-crossing energy transfer and not an event rate, and the step
between them remains blocked by the cancellation measured in
Sec.~\ref{sec:pf4b}.  The constant multiplying the adiabaticity is not
claimed, and the exploratory observation that a first-order singularity
gives close to minus two while a second-order pole gives about
$-1.1$ is not sealed.  Nothing in this section is a claim about quantum
field theory or about any physical process.  The campaign also stopped
here deliberately.  Two runs had already corrected two declaration
errors, and a third refinement aimed at a constant that was recognised
after measurement rather than predicted before it would have been
fitting the analysis to a wanted answer.

\section{Five vacuous runs, four bars that failed for the declaration,
and one invalid record}
\label{sec:spine}

This arc contains the densest failure record in the campaign.  Each
entry below is a recorded fact with its design error named.

Five runs passed or reported on manifests containing none of the events
their claims were about.  The pulse-train registration of
Sec.~\ref{sec:pf4b} sealed itself unevaluable because every cell capped
without a reversal, and the direct-observable probe then measured the
family empty at all twelve probed cells.  The first conservation gate
passed all four of its bars on a product grid whose eight cells
recorded zero reversing trajectories, so its claim about the accounting
of apparent pair events was untested by its own pass.  The first
covariance gate ran its audits on four deterministic members that
produce no folds, and only its own anti-vacuity bar caught it.  The
first ceiling run computed its classical oscillation statistic on a
curve that was exactly constant, because the declared field reversed
every member at every delay.  And the first summit run compared three
widths at different adiabaticities, so its cross-profile ratio test was
contaminated and could not be read.

The rule that came out of them tightened three times and is now
unconditional.  After the empty family it required a manifest to be
placed from a probe of the same observable at identical constants.
After the event-free grid it applied to gate runs as well as to summit
runs.  After the fold-free members it required a probe verifying event
presence for every bound cell and every bound member.  And after the
saturated classical arm, which was a new construction rather than an
inherited manifest and so slipped past the rule's earlier wording, it
became a requirement on any arm of any experiment.  Every registration
in this campaign must now cite a committed probe showing that its
manifest contains the variation its claim is about.

Four further bars failed for reasons belonging to the declaration
rather than to the model, and each is a different mistake.  The first
sealed axis comparison set its bars in binomial units, which tested a
model nobody had claimed, and the correction declared in writing that
its bars would be ratios to the model's own training inadequacy before
that registration existed.  The first ceiling bar of $0.15$ sat below
its own statistic's noise floor, which a later null measurement put at
a mean of $0.2165$ on a curve of that length, so no curve could have
met it.  The gate on the Lorentzian profile set a timestep bar of
$10^{-8}$ that the declared timestep could not deliver, missing at
$3 \times 10^{-7}$ for a reason with nothing to do with the profile.
And a conservation bar of $10^{-10}$ was missed by a factor of two,
its declared repair of halving the timestep was refuted by the rerun,
and the drift turned out to be a floating-point inconsistency in the
tilt, so the bar was eventually met by fixing an evaluation rather than
by integrating harder.

Two more entries belong in the same record.  A bar requiring a
coefficient of determination of $0.99$ was declared over two profiles
when only one of them was predicted to be exactly exponential, and the
other missed it at $0.9886$, which is itself informative about the
order of its singularity.  And two runners in the summit sequence
imported a fitting helper that regresses against the reciprocal of the
adiabaticity rather than the adiabaticity itself, which produced
positive exponents for a quantity already committed as negative.  The
error was found by comparison against a committed record, the runners
were corrected, and the declared protocols were rerun with no bar and
no design changed.

One of those two runners had already written its record into the
repository, and what followed is a failure of the record keeping rather
than of a declaration.  The record committed for the span ladder of
Sec.~\ref{sec:summit} carries the positive exponents the defect
produces, so its exponent items are invalid, and it stays in the
repository unedited.  The corrected rerun was written over the same
path in the working tree and was never committed, so for a time the
repository held an invalid record and nothing reading the repository
could see the correction.  The correction is now committed at its own
path, which names the record it supersedes, and the runner writes to
that new path so that a rerun cannot overwrite a committed record.  The
span-ladder convergence numbers do not pass through the fit and are
identical in both records, so the diagnosis of algebraic tails stands.
The two exponent numbers move, the worst relative exponent change from
$0.28$ to $0.23$ percent and the spread at the furthest span from
$0.356$ to $0.359$ percent, both far inside their bars, and the verdict
does not change.  Every exponent number reported for that run in
Sec.~\ref{sec:summit} and in Table~\ref{tab:ledger3} is read from the
corrected record.

Five of these failures would have been reported as successes by a
pipeline that did not check whether its manifest contained its events,
and the invalid record would have been read as current by anything that
trusted the repository.

\section{What this paper establishes}
\label{sec:claims}

\begin{table*}[t]
\caption{The ledger of the instrument layer and the sealed Schwinger
negative.  Verdicts are recorded exactly as computed from the declared
items of each run.  Proved means a statement about smooth maps under
its stated hypotheses.  Measured in model means an exact measurement in
a declared finite construction with append-only evidence records and
bars, grids, seeds, and estimators fixed before the run.  Key numbers
are quoted from the committed records.  Tables~\ref{tab:ledger2} and
\ref{tab:ledger3} carry the gates, the ceiling, the decay gate, the
summit sequence, and the scope of the whole arc.}
\label{tab:ledger}
\begin{ruledtabular}
\footnotesize
\begin{tabular}{L{2.5cm}lL{9.4cm}L{3.0cm}}
Run & Verdict & Key measurements & Label \\
\colrule
Fold kinematics & proved & a nondegenerate critical point of the
observed time changes the unsigned branch count by exactly two and the
signed count by exactly zero, under the stated smoothness and
genericity hypotheses & proved \\
PF0-FREEZE-001 instrument net & SEALED & branch counts 0, 1, 2 across
the creation fold with signed count 0 at every slice, the reverse order
across the annihilation fold, the cubic critical point classified
degenerate and never a fold, the double fold giving band counts 1, 3,
3, 3, 1 with signed count $+1$ and both critical points at
$\pm 1/\sqrt3$, separation exponent within $10^{-12}$ of one half,
cross-language branch locations agreeing to $4.5\times10^{-16}$,
classifier tolerances derived from a worst residual of
$4.6\times10^{-13}$ and a smallest fold curvature of 0.208 &
instrument \\
PF-1 structural stability & PASS & 500 trials at each of twelve
perturbation norms from $10^{-5}$ to 0.316, fold persistence exactly 1
through norm 0.123 and 0.956 at the largest with 22 of 500 lost, zero
degenerate classifications anywhere, median separation exponent 0.5
throughout and within $2.8\times10^{-6}$ of it at the largest
perturbation & measured in model, exploratory \\
PF-2 negative control & measured & a single member gives exactly 12
folds with 7 positive and 6 negative orientation segments at relative
drift $1.3\times10^{-7}$, and a 1000-member perturbed ensemble gives
exactly 12 folds in every member, 12000 in total with zero dispersion,
at mean relative drift $6.7\times10^{-10}$, with the same scenario
reproducing its outcome hash bit for bit on a second substrate &
measured in model, exploratory \\
PF-3 prior-art replication & replicated & the published closed form
agrees with the linear system it came from to $5.7\times10^{-14}$ and
with its own published time component to $1.4\times10^{-14}$, the
reversal threshold sits exactly at coupling 2, 27 sweep cells have
spacelike outgoing velocity and are recorded, the smoothed worldline
has zero observed-time crossings below threshold and exactly one above,
classified an annihilation fold & replicated \\
PF-4 pilot & measured & 18 cells at $5\times10^4$ members with drift
below $7\times10^{-7}$ in every cell, a transition from
$2.8\times10^{-4}$ at field 0.65 to 0.9115 at field 0.80 at gap 0.5,
only 2 of 18 cells in the measurable band, so every later manifest in
this family is probe-placed & measured in model, exploratory \\
PREREG-PF4-001 axis discrimination & neither-axis & all 20 probe-placed
cells measurable, both prescription controls agreeing per member at
$z=0.00$, zero-field null exactly zero, fitted effective-gap slope
41.24, the effective-gap model beating the Schwinger-shaped model by
3.28 on held-out error, but held-out weighted error 81.22 against a bar
of 4 because the bars were written in binomial units & measured in
model, exploratory \\
PREREG-PF4-002 bars against model error & PASS & held-out weighted
error 0.3842 of training error against a bar of 2, the Schwinger-shaped
model 3.1398 times worse against the same bar, fitted slope 43.18, ten
of twelve training and seven of eight held-out cells usable, controls
at $z=0.00$, zero-field null exactly zero & demonstrated in model \\
\end{tabular}
\end{ruledtabular}
\end{table*}

\begin{table*}[p]
\caption{The ledger of the empty family, the two mandatory gates, the
ceiling sequence, and the decay gate.  Both gates are recorded twice
because both first runs met every bar on a manifest containing none of
the events their claims were about.  The ceiling sequence is recorded
in full because four of its six runs corrected an error of the
declaration.  Key numbers are quoted from the committed records.}
\label{tab:ledger2}
\begin{ruledtabular}
\scriptsize
\begin{tabular}{L{2.5cm}L{1.9cm}L{9.0cm}L{2.7cm}}
Run & Verdict & Key measurements & Label \\
\colrule
PREREG-PF4-003 pulse train & unevaluable & every one of four training
and three held-out cells capped at 4000 slabs without a reversal, and
every cell of the mandatory planted control likewise, so no fit exists
and the record carries a manifest-design failure & measured in model,
exploratory \\
PF4-004 direct probe & empty & zero reversals at all twelve probed
cells, three fields crossed with four velocities at identical
constants, at a cap of 300 slabs, so the successor registration cannot
bind under its own verification clause & measured in model,
exploratory \\
PF4 mechanism study & flag False & net drift per slab 3.3, 0.9, and 5.0
percent of the mean deposit magnitude, so 95 to 99 percent of every
deposit is cancelled, sign flips on 54, 68, and 27 percent of
consecutive deposits, and the third cell's drift is $+0.0096$ per slab,
carrying the event away from the fold & measured in model,
exploratory \\
PREREG-PF5-001 conservation & PASS, content untested & all 8 cells
evaluable, every census summing exactly to 20000 with zero missing,
largest relative energy residual $6.5\times10^{-7}$ against a bar of
$10^{-5}$, path-degree rule holding at every generic level on all 96
declared members with zero failures and zero refused levels, and zero
reversing trajectories in every cell, all 160000 members transmitted &
measured in model, exploratory \\
PREREG-PF5-002 probe-placed & PASS & 1998, 1952, 1881, and 2237
reversing members out of 20000 each, 8068 audited apparent pair events,
every census exact with zero missing, largest relative residual
$7.3\times10^{-7}$ against $10^{-5}$, zero path-degree failures and
zero refused levels on all 48 declared members & demonstrated in
model \\
PREREG-PF6-001 covariance & vacuous & every declared deterministic
member produced zero folds, so all four cells were audited for the
covariance of events that were not present, and the run's own
anti-vacuity bar recorded it vacuous rather than passing & measured in
model, exploratory \\
PREREG-PF6-002 probe-placed & PASS & fold counts 4, 2, 2, 2 matching
the member probe, 10 folds in total and every count even, first-fold
worldpoints equivariant under the declared shifts to
$3.1\times10^{-14}$, reparametrized rate ratio deviating from the
declared factor by exactly 0, boost commutation
$4.5\times10^{-13}$ and invariant drift $7.1\times10^{-11}$ with a
positive future-cone margin, and the observer count equal to the
detector-model prediction at all six declared observers on every member
& demonstrated in model \\
PF-7 ceiling, first attempt & FAIL & single-passage control
$5.06\times10^{-3}$ from the closed form against a bar of
$2\times10^{-3}$, quantum statistic 0.2238 against 0.30 because a grid
uniform in delay undersamples a chirp quadratic in delay, norm
conserved to $5.5\times10^{-14}$, and the classical arm returning
exactly 1.0 at every one of 24 delays so its statistic was computed on
a constant curve & measured in model, exploratory \\
PF-7b corrected grid and bound field & PASS, ceiling bar failed &
single-passage control within $1.49\times10^{-3}$, quantum statistic
0.8702 under the grid uniform in the squared half-delay, norm
$5.9\times10^{-14}$, classical fractions varying from 0.4645 to 0.5245,
but the classical statistic 0.3626 against a declared 0.15 and the
ratio 2.40 against 3 & measured in model, exploratory \\
PF-7c noise floor & PASS & the null floor measured rather than
asserted, null mean 0.2165 and ninety-fifth percentile 0.3310 over
20000 draws, so the declared 0.15 was unattainable, classical p-value
0.0262 and quantum p-value 0.0, and the ceiling reading withheld and
escalated exactly as the protocol required & measured in model,
exploratory \\
PF-7d ringing test & PASS, mechanism refuted & structure survives four
times the ensemble at p-value 0.02605 against 0.0262 at a quarter of
the members, but the spectral peak sits at period 12.52 against a
declared ringing period of 5.236, which is the lowest nonzero bin and
therefore a trend, the curve falling from 0.4942 to 0.4829 across the
range & measured in model, exploratory \\
PF-7e linear detrend & FAIL & removing a declared linear trend raised
the classical statistic from 0.4480 to 0.5014 and lowered its p-value
from 0.02605 to 0.01185, refuting the diagnosis, while the trend is
real at 4.19 standard errors and the quantum statistic is untouched at
0.8702 & measured in model, exploratory \\
PF-7f declared smooth model & PASS & removing an asymptote plus an
exponential approach at the declared pulse length drops the classical
statistic from 0.4480 to 0.2122 at a p-value of 0.757, the quantum
statistic measures 0.8724 after the same treatment, and the trend
remains real at 4.19 standard errors & measured in model,
exploratory \\
PF-8 decay gate & PASS & 441 generic levels on four members and 400
thermal polylines, complete-observer counts only 1 and 3 with every
change minus two, zero, or plus two, zero odd events and zero parity
failures, zero path-degree failures on 441 deterministic and 2000
thermal checks with signed count exactly $+1$ and spread 0, the blind
consumer hiding 147 of 533 crossings, flipping parity on 129 of 441
observations, producing 63 apparent decays, and matching its own
detector model on all 441 with zero mismatches, and lifetimes with
coefficient of variation 0.159 against the 1 a memoryless process
requires & measured in model, exploratory \\
\end{tabular}
\end{ruledtabular}
\end{table*}

\begin{table*}[p]
\caption{The ledger of the summit sequence and the scope of the whole
arc.  The exploratory runs are kept because three of them failed and
each failure named the error that the next run corrected.  The last row
is a non-claim and governs every row of this table and of
Tables~\ref{tab:ledger} and \ref{tab:ledger2}.}
\label{tab:ledger3}
\begin{ruledtabular}
\footnotesize
\begin{tabular}{L{2.5cm}lL{9.4cm}L{3.0cm}}
Run & Verdict & Key measurements & Label \\
\colrule
PF4-005 analyticity, velocity grid & FAIL & timestep control
$6.9\times10^{-11}$, the sech-squared cell at width three replicating
the earlier probe's slope at $-7.8154$ against $-7.8186$, Lorentzian
coefficients of determination 0.99997, 0.99991, and 0.99933 with
exponent ratios 1.0095, 1.0344, and 1.0658 against sech-squared ratios
0.5430, 0.6910, and 0.8098, but the velocity grid compared three widths
at different adiabaticities & measured in model, exploratory \\
PF4-005b adiabaticity grid & FAIL & Lorentzian exponents $-2.0661$,
$-2.0639$, and $-2.0607$, a spread of 0.26 percent, 3.3 percent above
the parameter-free value, at coefficients of determination 0.9998,
sech-squared exponents $-1.0490$, $-1.0478$, and $-1.0462$, a spread of
0.27 percent, refuting the declared expectation, timestep control
$1.8\times10^{-13}$, and the failed bar a coefficient of determination
of 0.99 declared over both profiles when only one was predicted to be
exactly exponential & measured in model, exploratory \\
PF4-006 gate content & FAIL & conserved quantity drifting at
$8.6\times10^{-14}$, complete census with every crossing transmitted,
translation covariance $8.8\times10^{-14}$, but moving the entry point
from twenty widths to thirty changing the transfer by 1.289, 1.295, and
1.291 percent against a bar of $10^{-8}$ because the tilt approaches
its asymptote algebraically, and a timestep bar the declared timestep
could not deliver, missing at about $3\times10^{-7}$ & measured in
model, exploratory \\
PF4-007 span ladder & PASS & Lorentzian successive relative changes
1.278, 0.586, 0.130, and 0.015 percent along the declared span ladder
while the sech-squared transfer changes by at most $1.1\times10^{-13}$
against a declared bar of $10^{-9}$,
the exponent changing by at most 0.23 percent between the nearest and
furthest spans, and width universality surviving at the furthest span
with a spread of 0.36 percent, all read from the corrected record
\texttt{pf4-007b} that supersedes the first one committed & measured in
model, exploratory \\
PF4-008 pole order & FAIL & the entry clause the Lorentzian failed at
$1.3\times10^{-2}$ passed at $5.7\times10^{-11}$, Gudermannian
exponents $-1.9333$, $-1.9286$, and $-1.9217$ at a spread of 0.61
percent against sech-squared exponents $-1.1028$, $-1.1015$, and
$-1.0997$ on the identical grid, a ratio of 1.75 refuting the
distance-only hypothesis, but conservation drifting at
$1.9\times10^{-10}$ against a bar of $10^{-10}$ & measured in model,
exploratory \\
PF4-008b halved timestep & FAIL & the declared repair refuted itself,
the drift measuring $2.0\times10^{-10}$ at half the step, so the drift
was not integration error and the repair was wrong & measured in model,
exploratory \\
PF4-008c corrected tilt evaluation & PASS & evaluating the tilt so that
it saturates together with its derivative dropped the worst drift to
$2.6\times10^{-12}$ and moved the exponents only in their seventh
digit, with entry-point independence $3.0\times10^{-9}$, timestep
independence $6.6\times10^{-13}$, and translation covariance
$2.5\times10^{-15}$, and all eight items passing & measured in model,
exploratory \\
PREREG-PF4-009 sealed exponent & PASS & four fitted exponents
$-1.929735$, $-1.924383$, $-1.916823$, and $-1.909268$ across widths
spanning a factor of 2.56, a spread of 1.07 percent against a declared
bar of two percent, coefficients of determination 0.9989 or better in
every cell, Hamiltonian drift at most $3.7\times10^{-13}$ against
$10^{-11}$, entry-point independence $9.2\times10^{-9}$ against
$10^{-7}$, every transfer inside the declared window, and no cell
nonfinite and none folding & demonstrated in model \\
Scope & non-claim & the kinematic statement is a theorem and everything
else here is measured in declared finite models at declared constants
and a declared step, the sealed observable is a per-crossing energy
transfer and not an event rate, the constant multiplying the
adiabaticity is not claimed, apparent decay here is a measurement of
consumer blindness and never a property of the worldline, and no result
in this paper is evidence that physical pair creation or physical
particle decay works this way & non-claim \\
\end{tabular}
\end{ruledtabular}
\end{table*}

Tables~\ref{tab:ledger}, \ref{tab:ledger2}, and \ref{tab:ledger3}
register every run of the arc with its recorded verdict.  Read together
they answer the founding question inside its declared scope, and the
answer has four parts.

The kinematic part is settled and is the only part that is a theorem.
A fold of the observed time makes exactly two branches of opposite
orientation and leaves the signed count alone.  The instrument that
finds them was validated against exact controls before any dynamics
ran, its tolerances were derived from measured margins rather than
chosen, and the folds it finds survive perturbation.

The dynamical part is a negative in the one family the campaign could
populate and seal.  Reversal suppression there is a Gaussian tail of
the initial-condition measure in a measured effective gap, and the
Schwinger-shaped variable does not compete.  The family constructed to
escape that verdict by removing the measure entirely turned out to
contain no events, because its alternating slabs cancel almost all of
every deposit and the residual has no fixed sign.  Both mandatory gates
hold on cells that contain events, conservation with every census exact
and the orientation charge rule never failing, and covariance with
every count invariant under the declared passive transformations and
every observer-dependent count predicted exactly by a declared detector
model.

The ceiling part is a measured limit.  The model reproduces the
kinematics of pair events and the phenomenology of apparent decay under
restricted observation, and it does not reproduce quantum phase
structure.  A two-level mode interferes at a level no noise floor
approaches while the matched classical ensemble carries only the smooth
delay dependence its own pulse overlap predicts.  This caps the model
as a classical representation and does not touch the kinematic
statement or any other measured result here.

The fourth part is the one sealed positive and it is narrow.  There
exists a constructible family whose suppression exponent is fixed by
the distance to the field's nearest complex singularity, independent of
the profile's width across a factor of 2.56 to within 1.07 percent,
measured on deterministic crossings where no distribution exists whose
tail could produce it.  That answers the construction question the
sealed negative left open, for the per-crossing transfer and for
nothing else.  The step from a transfer to an event rate is blocked by
the cancellation the mechanism study measured, and no claim about a
pair production rate follows.

Three objects here do not depend on the substrates measured.  The first
is the parity argument of Sec.~\ref{sec:pf8}, which excludes a
same-species decay of one particle into two for complete continuous
worldlines under generic projections by arithmetic rather than by
search.  The second is the lawful-observer clause, which allows
observer-dependent particle number exactly when a declared detector
model predicts every count, and which held on 441 observations with
zero mismatches for a consumer that was blind by declaration.  The
third is the standing rule of Sec.~\ref{sec:spine}, which five failures
in this arc paid for.

Nothing above is evidence that physical spacetime is a projection, that
physical pair creation proceeds by folding, or that physical particle
decay is a measurement artifact.  Every dynamical statement here is
about a declared finite model at declared constants and a declared
step.

\section{Methods and reproducibility}
\label{sec:methods}

All runs follow the campaign's standing discipline.  Model families,
constants, grids, seeds, integrators, tolerances, estimators, stopping
rules, and falsification bars were declared in the track documents
before the runs they govern, and claim-bearing runs are governed by a
preregistration whose file hash is recorded in an append-only seal
ledger at its sealing commit, verifiable against the repository's git
history.  The nine sealed documents of this arc are the tolerance
freeze, three registrations on the Schwinger question, two on
conservation, two on covariance, and the sealed analyticity exponent.
Verdicts are computed from the declared items by the governed runner
rather than written afterward.

Every result-bearing run writes an append-only JSON evidence record
containing the declared parameters, the measured outcomes, the computed
verdict, the code commit, the dependency versions, and a canonical
SHA-256 hash of the record body.  The records for this paper are the
instrument-net and ensemble records, the structural-stability pilot,
the replication record, the pilot and the three sealed records of the
Schwinger question, the direct probe and the mechanism study of the
pulse-train family, the two conservation instrument and gate records
with their placement probe, the two covariance instrument and gate
records with their member probe, the six ceiling records with their two
field probes, the decay gate with its blind probe, the nine records of
the summit sequence, and the sealed analyticity record, all in
\texttt{results/} in the campaign repository, each with its committed
runner beside it.  Failed and vacuous records are retained unchanged
and each corrected declaration is a separate record, so the sequence of
verdicts is itself part of the evidence.  A reducer may summarize
records but may not rewrite them, and a record whose runner was later
found defective is superseded by a corrected record written to a new
path rather than overwritten, which is why the span ladder of
Sec.~\ref{sec:summit} has two records and the paper reads the second.

Every measured number and every figure coordinate in this paper is read
from those committed records.  The figure tables are generated
deterministically by a committed script,
\texttt{make\_\allowbreak figdata\_\allowbreak pf.py}, which uses only
the Python standard library, has no free parameters, and asserts each
value typed in this paper against the record value it came from before
it writes anything, so a record that changed would break the build
rather than disagree quietly with the text.  Declared parameters quoted
in the text, including the instrument-net controls, the frozen
classifier tolerances, the seal ledger entries, and the bars of every
registration, are checked the same way against the track documents and
preregistrations that declared them.

The instrument layer runs in two independent implementations, one in
MATLAB with event-located adaptive integration and one in Python with
fixed-step symplectic integration, and the frozen tolerances bind both
paths separately, with an adaptive path at relative tolerance
$10^{-10}$ and a fixed-step path at step $10^{-3}$ valid to evolution
parameter 40 \cite{HairerLubichWanner}.  Cross-substrate determinism is
part of the freeze, and the same scenario on a workstation and in a
container reproduces both its scenario hash and its outcome hash bit
for bit.  The Sauter-family ensembles are brute-force with no
importance sampling and no adaptive seeding, with seeds derived from
declared formulas in the governed runners, and every trajectory is
counted in every denominator including the numerically failed, the
capped, and the transmitted.  The summit sequence integrates single
deterministic crossings by fourth-order Runge-Kutta at declared steps
of $2\times10^{-4}$ and $10^{-4}$, with the extended
two-degree-of-freedom Hamiltonian's drift, the entry-point
independence, the timestep independence, and the translation covariance
all measured per run rather than assumed.  All runs executed on a
single workstation and require no special hardware.

\begin{thebibliography}{99}

\bibitem{Whitney1955}
H.~Whitney,
On singularities of mappings of Euclidean spaces I. Mappings of the
plane into the plane,
Ann. Math. \textbf{62}, 374 (1955).

\bibitem{ArnoldGZ}
V.~I.~Arnold, S.~M.~Gusein-Zade, and A.~N.~Varchenko,
\emph{Singularities of Differentiable Maps, Volume 1}
(Birkh\"auser, Boston, 1985).

\bibitem{Stueckelberg1941}
E.~C.~G.~Stueckelberg,
La signification du temps propre en m\'ecanique ondulatoire,
Helv. Phys. Acta \textbf{14}, 588 (1941).

\bibitem{Feynman1949}
R.~P.~Feynman,
The theory of positrons,
Phys. Rev. \textbf{76}, 749 (1949).

\bibitem{HorwitzPiron1973}
L.~P.~Horwitz and C.~Piron,
Relativistic dynamics,
Helv. Phys. Acta \textbf{46}, 316 (1973).

\bibitem{Land2016}
M.~Land,
Pair production in classical Stueckelberg-Horwitz-Piron
electrodynamics,
J. Phys.: Conf. Ser. \textbf{845}, 012025 (2017),
arXiv:1604.01625.

\bibitem{Sauter1931}
F.~Sauter,
\"Uber das Verhalten eines Elektrons im homogenen elektrischen Feld
nach der relativistischen Theorie Diracs,
Z. Phys. \textbf{69}, 742 (1931).

\bibitem{HeisenbergEuler1936}
W.~Heisenberg and H.~Euler,
Folgerungen aus der Diracschen Theorie des Positrons,
Z. Phys. \textbf{98}, 714 (1936).

\bibitem{Schwinger1951}
J.~Schwinger,
On gauge invariance and vacuum polarization,
Phys. Rev. \textbf{82}, 664 (1951).

\bibitem{Dunne2004}
G.~V.~Dunne,
Heisenberg-Euler effective Lagrangians: Basics and extensions,
in \emph{From Fields to Strings: Circumnavigating Theoretical Physics}
(World Scientific, Singapore, 2005),
arXiv:hep-th/0406216.

\bibitem{Landau1932}
L.~D.~Landau,
Zur Theorie der Energie\"ubertragung II,
Phys. Z. Sowjetunion \textbf{2}, 46 (1932).

\bibitem{Zener1932}
C.~Zener,
Non-adiabatic crossing of energy levels,
Proc. R. Soc. London A \textbf{137}, 696 (1932).

\bibitem{Stueckelberg1932}
E.~C.~G.~Stueckelberg,
Theorie der unelastischen St\"osse zwischen Atomen,
Helv. Phys. Acta \textbf{5}, 369 (1932).

\bibitem{Shevchenko2010}
S.~N.~Shevchenko, S.~Ashhab, and F.~Nori,
Landau-Zener-St\"uckelberg interferometry,
Phys. Rep. \textbf{492}, 1 (2010).

\bibitem{Dykhne1962}
A.~M.~Dykhne,
Adiabatic perturbation of discrete spectrum states,
Sov. Phys. JETP \textbf{14}, 941 (1962).

\bibitem{DavisPechukas1976}
J.~P.~Davis and P.~Pechukas,
Nonadiabatic transitions induced by a time-dependent Hamiltonian in the
semiclassical/adiabatic limit: The two-state case,
J. Chem. Phys. \textbf{64}, 3129 (1976).

\bibitem{HairerLubichWanner}
E.~Hairer, C.~Lubich, and G.~Wanner,
\emph{Geometric Numerical Integration}, 2nd ed.
(Springer, Berlin, 2006).

\bibitem{ZenodoEntropy}
A.~H.~Bond,
Entropy across singular projections: Fiber multiplicity, caustics, and
observer-relative information,
Zenodo (2026), doi:10.5281/zenodo.21789011.

\bibitem{ZenodoCayley}
A.~H.~Bond,
The pair-creation threshold of impulsive Stueckelberg-Horwitz-Piron
scattering is a Cayley-transform pole,
Zenodo (2026), doi:10.5281/zenodo.21790096.

\bibitem{ZenodoSchwinger}
A.~H.~Bond,
A sealed negative for Schwinger scaling in a resolved-field
hidden-dynamics family,
Zenodo (2026), doi:10.5281/zenodo.21798545.

\bibitem{ZenodoGames}
A.~H.~Bond,
Games and decisions under projection: Blackwell wedges, budget flips,
and prospect signatures from declared consumers,
Zenodo (2026), doi:10.5281/zenodo.21818268.

\bibitem{ZenodoGen2}
A.~H.~Bond,
Five audits under one consumer discipline: Renormalization, quantum
Darwinism, lattice hydrodynamics, fluctuation theorems, and gauge
redundancy in exact finite models,
Zenodo (2026), doi:10.5281/zenodo.21821378.

\bibitem{ZenodoCrypto}
A.~H.~Bond,
Engineered observation limits in exact finite models,
Zenodo (2026), doi:10.5281/zenodo.21833690.

\end{thebibliography}

\end{document}
